%============================================================ 4. 제안 방법
\section{제안 방법}
\label{sec:method}

\subsection{2계층 비동기 구조}
\label{sec:arch}

제안 체계는 방어 판단을 전략 계층과 기동 계층으로 분리한다(Fig.~\ref{fig:architecture}).
전략 계층은 \emph{어느 방어정이 어느 적 무리를 맡는가}를 정하고, 기동 계층은 \emph{배정된
무리를 막기 위해 해면 어디에 그물벽을 놓는가}를 정한다. 두 계층은 같은 주기로 결정을 내리지만
\textbf{응답 속도가 다르다}. 기동 계층의 순전파는 순간에 끝나는 반면, 전략 계층의
언어모델 호출은 응답에 지연이 발생한다(\ref{sec:res_latency}절). 느린 전략 계층 호출만
비동기로 분리하여, 호출이 진행되는 동안 기동 계층은 직전 배정으로 결정을 계속하고 응답이
도착하면 다음 결정부터 새 배정을 반영한다. 이하에서 쓰는 기호는
Table~\ref{tab:nomenclature}에 모았다.

\subsection{적 무리 군집화}
\label{sec:cluster}
% [H]: 절 제목 바로 아래 고정(float 패키지). 단 폭 그림이라 같은 단에 들어간다.
\begin{figure}[H]
\centering
\includegraphics[width=\columnwidth]{fig5.png}
\caption{방위각 정렬, 군집화, 군집 배정}
\label{fig:clustering}
\end{figure}

두 판단 계층은 개별 적선이 아니라 \emph{무리}를 단위로 배정한다. 따라서 매 결정마다 생존
적선을 모선 기준 방위로 묶어 최대 $C$ 개의 무리로 요약한다. 절차는 Fig.~\ref{fig:clustering}의
네 단계로 이루어진다. 모선 기준 방위각을 정렬하고, 임계 $\delta_{\mathrm{gap}}$ 를 넘는 방위
빈틈을 탐색하며, 빈틈에서 잘라 $1$--$C$개 무리를 확정한 뒤, 무리별 위협도와 요격점(무리
중심--모선 선분 위)을 산출한다. 적선 $m$ 의 위치를 $\mathbf{e}_m = (e_{m,x}, e_{m,y})$, 모선 위치를
$\mathbf{m} = (m_x, m_y)$ 라 하면, 적선의 모선 기준 방위 $\vartheta_m$ 은 항법 규약(북쪽
$0^\circ$, 시계 방향)으로 정의된다.
\begin{equation}
\vartheta_m = \operatorname{atan2}\!\big(e_{m,x} - m_x,\; e_{m,y} - m_y\big) \bmod 360^\circ
\label{eq:bearing}
\end{equation}

고정 방위 구간(bin)은 한 무리가 경계에 걸리면 둘로 쪼개지므로, \emph{방위 간격(gap)}으로
무리 수를 데이터에 맞추는 적응형 군집화를 쓴다. 생존 적선
$n = |\mathcal{A}|$ 척의 방위를 오름차순으로 정렬하여 $\vartheta_{(1)} \le \dots \le
\vartheta_{(n)}$ 이라 하고, 이웃한 방위 사이의 간격 $g_{(q)}$ 를 원형으로 정의한다.
\begin{equationfn}
g_{(q)} = \vartheta_{(q+1)} - \vartheta_{(q)}\ (q < n), \quad
g_{(n)} = \vartheta_{(1)} + 360^\circ - \vartheta_{(n)}
\label{eq:gap}
\end{equationfn}
마지막 항 $g_{(n)}$ 은 $\vartheta_{(n)}$ 에서 북쪽을 지나 $\vartheta_{(1)}$ 로 되돌아오는
이음새 간격으로, 북쪽에 걸친 무리가 쪼개지는 것을 막는다.
원 위에서는 임계 $\delta_{\mathrm{gap}}$ 을 넘는 간격 $k$ 개가 곧 무리 $k$ 개이므로 그 수를
활성 무리 수 $\hat{C}$ 로 삼되, $1$ 이상 $\min(C, n)$ 이하로 자른다.
\begin{equation}
\hat{C} = \mathrm{clip}\!\Big(\big|\{\, q : g_{(q)} > \delta_{\mathrm{gap}} \,\}\big|,\; 1,\; \min(C, n)\Big)
\label{eq:ncluster}
\end{equation}
가장 큰 간격 $\hat{C}$ 개를 무리 경계로 삼고, 그 가운데 최대 간격을 원형 순서의 이음새로
열어 이음새 다음 적선부터 방위 순으로 순회하며 경계를 지날 때마다 무리 번호를 하나씩
올린다. 이렇게 얻은 라벨 $j_m \in \{1, \dots, \hat{C}\}$ 로 무리 $j$ 의 구성원 집합
$\mathcal{A}_j = \{m : j_m = j\}$ 가 정해진다. 무리 $j$ 는 척수 $n_j$, 중심
$\mathbf{c}_j$, 각도 퍼짐 $\Delta\vartheta_j$, 모선 방향 접근 속도 $v^{\mathrm{app}}_j$ 의 네 가지
양으로 요약한다. 각도 퍼짐은 방위의 산술 편차가 $0^\circ/360^\circ$ 경계에서 깨지므로 방위 단위벡터
$\hat{\mathbf{b}}_m = (\cos\vartheta_m, \sin\vartheta_m)$ 의 평균 길이(원형 결집도)
$\bar{R}_j$ 로 잰다(한 점에 모이면 $0^\circ$). 접근 속도는 적선 침로 $\psi_m$ 의 단위벡터와
적선에서 모선을 향한 단위벡터 $\hat{\mathbf{d}}_m$ 의 내적에 적선 속력 $v_e$ 를 곱한 값이다(이탈
시 음수).
\begin{equation}
\begin{aligned}
n_j &= |\mathcal{A}_j|, \qquad
\mathbf{c}_j = \frac{1}{n_j} \sum_{m \in \mathcal{A}_j} \mathbf{e}_m,\\
\Delta\vartheta_j &= 90^\circ \sqrt{1 - \bar{R}_j}, \qquad
\bar{R}_j = \Big\lVert \frac{1}{n_j} \sum_{m \in \mathcal{A}_j} \hat{\mathbf{b}}_m \Big\rVert,\\[2pt]
v^{\mathrm{app}}_j &= \frac{v_e}{n_j} \sum_{m \in \mathcal{A}_j}
  \big\langle (\sin\psi_m, \cos\psi_m),\, \hat{\mathbf{d}}_m \big\rangle,\\
\hat{\mathbf{d}}_m &= (\mathbf{m} - \mathbf{e}_m) / \lVert \mathbf{m} - \mathbf{e}_m \rVert
\end{aligned}
\label{eq:clustersum}
\end{equation}
본 연구는 $C = 4$, $\delta_{\mathrm{gap}} \approx 12^\circ$ 를 쓴다.
% ★ Fig.7(fig:observation) 은 §4.4.1 이 아니라 여기(§4.2 끝, 5쪽 본문)에서 선언한다.
%   2단 조판의 figure* 는 선언된 쪽에 놓일 수 없어, §4.4.1 본문에서 선언하면 7쪽으로 밀린다.
%   여기서 선언하면 6쪽 상단 --- §4.4.1 이 시작하는 쪽 --- 에 놓인다(2026-09-16 배치 실측).
\begin{figure*}[!tp]
\centering
\includegraphics[width=0.62\textwidth]{fig6.png}
\caption{래스터화된 15채널 관측의 예(파상 포메이션, 통영 매물도 북측)}
\label{fig:observation}
\end{figure*}

이 규칙 아래에서 집중 포메이션은 무리 하나, 양동 포메이션은 세 방위의 무리 셋으로
요약된다. 파상 포메이션은 단(rank)마다 방위를 $18^\circ$ 씩 벌려 생성되므로
$\delta_{\mathrm{gap}}$ 을 넘어 대개 단마다 별개의 무리가 되며, 단 사이의 거리차는 무리
중심의 반경 차이와 \ref{sec:obs}절의 격자 관측에 함께 남는다. 척수와 중심은
\ref{sec:heuristic}절의 위협도와 요격점을, 네 요약량 전부는 \ref{sec:commander}절의
지휘관 입력을 구성한다.

\subsection{기하 휴리스틱 기준 방법}
\label{sec:heuristic}

\ref{sec:cluster}절의 무리 요약 위에서 같은 문제를 기하 규칙만으로 푸는 기준 방법으로,
\ref{sec:results}절의 비교 기준이자 \ref{sec:bc}절의 행동 복제 목표다.

\textbf{배정.} 무리 $j$ 의 위협도 $\tau_j$ 는 척수 $n_j$ 에 모선 근접도(무리 중심
$\mathbf{c}_j$ 에서 모선 $\mathbf{m}$ 까지의 거리를 해역 반폭 $R_w$ 로 나눈 값을 1에서 뺀 것)를
곱한 양이고, 요격점 $\mathbf{I}_j$ 는 무리 중심에서 모선을 향해 선분의 $\lambda_{\mathrm{I}}$
비율만큼 옮긴 점이다.
\begin{equation}
\begin{aligned}
\tau_j &= n_j \cdot \mathrm{clip}\!\left(1 - \frac{\lVert \mathbf{c}_j - \mathbf{m}\rVert}{R_w},
         \, 0.05,\, 1\right),\\
\mathbf{I}_j &= \mathbf{c}_j + \lambda_{\mathrm{I}}\,(\mathbf{m} - \mathbf{c}_j), \qquad \lambda_{\mathrm{I}} = 0.55
\end{aligned}
\label{eq:threat}
\end{equation}
위협도는 크고 가까운 무리일수록 크며, 하한 $0.05$ 는 해역 가장자리의 먼 무리도 척수에
비례하는 최소 위협을 갖게 한다. 요격점은 $\lambda_{\mathrm{I}} = 0$ 이면 무리 중심, $1$ 이면
모선이다. 방어정이 적보다 느리므로 무리 중심에 가까운 요격점은 적이 먼저 지나쳐
버린다. 모선 쪽으로 치우친 $\lambda_{\mathrm{I}} = 0.55$ 는 방어정이 적보다 먼저 도달할 수
있는 반경 안에 요격점을 두어 차단의 실현성을 확보하는 값이다.
위협 상위 $P$ 개 무리에 대해 탐욕적 1:1 배정을 한다. 방어정 $p$ 의 위치를 $\mathbf{a}_p$,
요격점까지의 거리를 $d_{pj} = \lVert \mathbf{a}_p - \mathbf{I}_j \rVert$, 직전 결정에서 $p$ 가
맡았던 무리를 $j_p^{\text{prev}}$, 지시함수를 $\mathbb{1}[\cdot]$ 이라 하면, 배정 비용
$\tilde{d}_{pj}$ 는 이 거리에서 직전 배정을 유지하는 쌍에만 보너스 $\beta = 6000\unit{m}$ 를 감한 값이다.
\begin{equation}
\tilde{d}_{pj} = d_{pj} - \beta \cdot \mathbb{1}[\,j = j_p^{\text{prev}}\,]
\label{eq:sticky}
\end{equation}
활성 무리가 $P$ 보다 적으면 남는 방어정은 부하가 가장 큰 무리에 안쪽 층($\lambda_{\mathrm{L}}
= 0.18$, 방위 $\pm 40^\circ$)으로 덧붙인다.

\textbf{경유점.} 그물벽은 배정 무리의 요격점 $\mathbf{I}$ 를 가운데 두고 접근 방향에 수직인
길이 $L_{\mathrm{net}} = 450\unit{m}$ 의 선분이므로, 경유점 $\mathbf{t}_k$ 는 그 양 끝 $K = 2$
개다. 모선에서 요격점을 향한 단위 벡터를 $\hat{\mathbf{r}}$, 그에 수직인 단위 벡터를
$\hat{\mathbf{r}}^{\perp}$ 라 하면
\begin{equationsmall}
\mathbf{t}_k = \mathbf{I} + \left(k - \tfrac{K-1}{2}\right) L_{\mathrm{net}}\,
               \hat{\mathbf{r}}^{\perp}, \qquad k = 0,\dots,K-1
\label{eq:heur_target}
\end{equationsmall}
이다. 각 $\mathbf{t}_k$ 는 학습 정책과 같은 격자·유효 마스크·간격 제약(\ref{sec:validmask}절) 안의
최근접 픽셀로 스냅하므로 두 방법의 행동공간은 같다. 휴리스틱은 기하로 정해진 한 점만
지목하며, 다른 방어정의 벽, 파상 공격의 뒤 단, 지형의 길목은 고려하지 않는다. 반면 학습된
점수맵(\ref{sec:policy}절)은 이 요소들을 반영하여 해면 전체의 후보 위치를 한꺼번에 비교한다.
따라서 \ref{sec:res_maneuver}절의 차이는 표현력이 아니라 \emph{결정 규칙}의 차이다.

\subsection{점수맵 기반 기동 정책}
\label{sec:policy}

\subsubsection{다채널 격자 관측}
\label{sec:obs}

기동 계층의 입력은 좌표 목록이 아니라, 모선을 중앙에 두고 관측 반폭 $E$ 의 전장을
$n_{\mathrm{g}} \times n_{\mathrm{g}}$ 격자로 나눈 다채널 지도이다. 총 15개 채널은 모든
방어정이 공유하는 전역 9채널, 방어정마다 다른 3채널, CoordConv~\citep{liu2018coordconv}
좌표 3채널로 구성되며(Fig.~\ref{fig:tensorstack}; 채널별 실제 화면은
Fig.~\ref{fig:observation}), 각 채널의 정의와 값 범위를 Table~\ref{tab:channels}에
정리하였다.

\begin{figure}[!htbp]
\centering
\includegraphics[width=\columnwidth]{fig7.png}
\caption{15채널 관측의 구성: 전역·방어정별·좌표 채널의 격자 중첩과 격자 밖 자체 상태 벡터}
\label{fig:tensorstack}
\end{figure}

\begin{table*}[!tp]
\centering
\caption{15채널 격자 관측의 구성. 전역 채널은 월드당 한 번 계산하여 모든 방어정이 공유하고,
         방어정별 채널만 방어정마다 다르다. 정의 열의 식 번호는 본문의 수식이다.}
\label{tab:channels}
\footnotesize
\renewcommand{\arraystretch}{1.15}
\begin{tabular}{@{}clllc@{}}
\toprule
\# & 묶음 & 채널 & 정의 & 범위 \\
\midrule
1  & \multirow{9}{*}{전역}   & 적 밀도        & 픽셀 내 생존 적선 수$/M$, 식~\eqref{eq:ch_enemy}      & $[0,1]$ \\
2  &                         & 적 진행 $x$    & 픽셀 내 적선 헤딩 단위벡터의 평균, 식~\eqref{eq:ch_enemy} & $[-1,1]$ \\
3  &                         & 적 진행 $y$    & 위와 같음                                              & $[-1,1]$ \\
4  &                         & 적 위협도      & 모선 도달 임박도의 픽셀 내 최대, 식~\eqref{eq:ch_threat} & $[0,1]$ \\
5  &                         & 아군 위치      & 픽셀 내 생존 방어정 수$/P$                             & $[0,1]$ \\
6  &                         & 설치된 그물    & 물리 격자의 그물 점유를 최대값 풀링                     & $\{0,1\}$ \\
7  &                         & 원점           & 모선 반경 안의 픽셀                                    & $\{0,1\}$ \\
8  &                         & 요격 구역      & $r_{\min} \le \lVert \mathbf{x}_i - \mathbf{m} \rVert \le r_{\max}$ 인 픽셀 & $\{0,1\}$ \\
9  &                         & 육지           & 고해상도 육지 마스크의 면적비  & $[0,1]$ \\
\midrule
10 & \multirow{3}{*}{방어정별} & 자기 위치     & 방어정 $p$ 가 놓인 픽셀                                & $\{0,1\}$ \\
11 &                         & 배정 요격점    & $\mathbf{I}_p$ 가 놓인 픽셀 (미배정이면 전부 0)        & $\{0,1\}$ \\
12 &                         & 유효 영역      & 유효 마스크 $\mathcal{V}_p$, 식~\eqref{eq:validset}    & $\{0,1\}$ \\
\midrule
13 & \multirow{3}{*}{좌표}   & $x$            & $(x_i - m_x)/E$, 식~\eqref{eq:ch_coord}                & $[-1,1]$ \\
14 &                         & $y$            & $(y_i - m_y)/E$                                        & $[-1,1]$ \\
15 &                         & $r$            & $\mathrm{clip}(\lVert \mathbf{x}_i - \mathbf{m} \rVert / E, 0, 1)$ & $[0,1]$ \\
\bottomrule
\end{tabular}
\end{table*}

\textbf{래스터화와 정규화.} 픽셀 크기는 $\ell = 2E/n_{\mathrm{g}}$ 이고, 월드 좌표
$\mathbf{x} = (x, y)$ 는 모선 위치 $\mathbf{m} = (m_x, m_y)$ 를 중심으로 다음 픽셀 인덱스
$i = (\iota_x, \iota_y)$ 로 사상된다. 각 축에 $E$ 를 더하는 것은 원점을 모선에서 격자의
모서리로 옮기기 위함이다.
\begin{equation}
\begin{aligned}
\iota_x &= \mathrm{clip}\!\left(\left\lfloor \frac{x - m_x + E}{\ell} \right\rfloor,\, 0,\, n_{\mathrm{g}}-1\right),\\
\iota_y &= \mathrm{clip}\!\left(\left\lfloor \frac{y - m_y + E}{\ell} \right\rfloor,\, 0,\, n_{\mathrm{g}}-1\right)
\end{aligned}
\label{eq:world2pix}
\end{equation}
적 관련 채널은 각 개체의 연속 위치를 식~(\ref{eq:world2pix})로 픽셀에 사상한 뒤 같은
픽셀에 속한 개체를 집계하여 값을 만들고, 값 범위가 $[0,1]$ 또는 $[-1,1]$ 에 놓이도록
정규화한다. 생존 적선 집합을 $\mathcal{A}$, 적선 $m$ 의
위치·헤딩을 $\mathbf{e}_m$, $\psi_m$, 위치가 픽셀 $i$ 로 사상되는 생존 적선의 집합을
$\mathcal{A}_i = \{m \in \mathcal{A} : \iota(\mathbf{e}_m) = i\}$, 그 수를 $N_i = |\mathcal{A}_i|$,
적선 속력 $v_e$ 로 에피소드 상한 $T_{\max}$ 동안 달릴 수 있는 최대 거리를 $v_e T_{\max}$ 라
하면, 픽셀 $i$ 의 적 밀도·진행 방향·위협도 채널 $c^{\mathrm{den}}_i$, $c^{\mathrm{vel}}_i$,
$c^{\mathrm{thr}}_i$ 는
\begin{equationsmall}
c^{\mathrm{den}}_i = \frac{N_i}{M}, \quad
c^{\mathrm{vel}}_i = \frac{1}{\max(N_i, 1)} \sum_{m \in \mathcal{A}_i} (\sin\psi_m, \cos\psi_m)
\label{eq:ch_enemy}
\end{equationsmall}
\begin{equation}
c^{\mathrm{thr}}_i = \max_{m \in \mathcal{A}_i}
   \left(1 - \mathrm{clip}\!\left(\frac{\lVert \mathbf{e}_m - \mathbf{m} \rVert}{v_e T_{\max}},\, 0,\, 1\right)\right)
\label{eq:ch_threat}
\end{equation}
이다. 아군 위치 채널은 같은 방식으로 $P$ 로 나누고, 설치된 그물
채널은 고해상도 물리 격자의 최대값 풀링, 육지 채널은 OpenStreetMap 해안선·섬 폴리곤을
래스터화한 실해역 마스크의 면적비 다운샘플이다(충돌 판정은 물리 격자에서 이진으로 한다).
좌표 채널은 픽셀 중심 $\mathbf{x}_i = (x_i, y_i)$ 에 대해
\begin{equationsmall}
(c^{x}_i,\, c^{y}_i) = \frac{\mathbf{x}_i - \mathbf{m}}{E}, \qquad
c^{r}_i = \mathrm{clip}\!\left(\frac{\lVert \mathbf{x}_i - \mathbf{m} \rVert}{E},\, 0,\, 1\right)
\label{eq:ch_coord}
\end{equationsmall}
로 두어 $[-1,1]$ 과 $[0,1]$ 에 놓는다. 격자에 담지 않는 방어정 $p$ 의 자체 상태 벡터
$\mathbf{o}_p$ 는 같은 반폭 $E$ 로 정규화한 위치 $\mathbf{a}_p$ 와 배정 요격점 $\mathbf{I}_p$,
헤딩 $\psi_p$ 의 단위벡터, 잔여 그물 수 $n_p$ 를 척당 보유량 $n_{\max}$ 로 나눈 비율, 전개 중
여부 $\delta_p \in \{0,1\}$ 로 구성되며, 배정이 없으면 요격점 성분은 0이다(지시함수
$\mathbb{1}[\text{배정}]$).
\begin{equation}
\begin{aligned}
\mathbf{o}_p = \Big[&\tfrac{1}{E}(\mathbf{a}_p - \mathbf{m});\; (\sin\psi_p, \cos\psi_p);\;
  \tfrac{n_p}{n_{\max}};\; \delta_p;\\
  &\mathbb{1}[\text{배정}]\,\tfrac{1}{E}(\mathbf{I}_p - \mathbf{m})\Big] \in \mathbb{R}^{8}
\end{aligned}
\label{eq:own}
\end{equation}

\textbf{불변성.} 길이는 $E$ 나 $\ell$ 로, 개수는 $M$·$P$·$n_{\max}$ 로, 시간은
$v_e T_{\max}$ 로 나뉘므로 모든 입력은 $[-1,1]$ 의 무차원 양이며, 격자 지도는 개체 수와
순서에도 불변이다. 따라서 해역 크기·격자 해상도·선박 수가 바뀌어도 입출력의 의미와 형태가
같고 통계 기반 정규화 층이 필요 없다($E$, $n_{\mathrm{g}}$ 등의 값은
Table~\ref{tab:scenario}). 합성곱은 평행이동 등변이라 절대 위치가 흐려지므로 좌표 채널을
함께 넣어 위치를 보존한다.

%-------------------------------------------------- 그림 4

\subsubsection{FiLM 조건화 U-Net과 점수맵}

\begin{figure*}[!tbp]
\centering
\includegraphics[width=0.75\textwidth]{fig8.png}
\caption{점수맵 정책의 구조: 방어정 상태 질의와 U-Net 특징맵의 내적으로 얻는 점수맵과 부화소 오프셋}
\label{fig:policy}
\end{figure*}

관측은 U-Net 구조~\citep{ronneberger2015unet}를 지나 입력과 같은 해상도의 특징 지도로
변환된다(Fig.~\ref{fig:policy}). 인코더 2단($50\to25\to13$)으로 수용 영역을 넓히고, 디코더
2단으로 원해상도를 복원하며, skip 연결이 인코더의 위치 정보를 디코더로 전달한다.

병목에서는 전역 평균 풀링으로 압축한 전황 벡터로부터 FiLM~\citep{perez2018film} 의
채널별 배율·이동 계수 $(\boldsymbol{\gamma}_{\mathrm{F}}, \boldsymbol{\beta}_{\mathrm{F}})$ 를
만들어 병목 특징을 변조한 뒤 디코더에 넘긴다. 병목 특징 지도를 $\mathbf{H}_b \in \mathbb{R}^{2w \times h \times h}$
($w = 32$, $h = 13$), 전역 평균 풀링을 $\mathrm{GAP}$, 계수를 내는 선형층의 가중치와 편향을
$\mathbf{W}_{\mathrm{F}}$, $\mathbf{b}_{\mathrm{F}}$, 채널별 곱을 $\odot$ 이라 하면 변조된 특징
$\tilde{\mathbf{H}}_b$ 는
\begin{equation}
\begin{aligned}
(\boldsymbol{\gamma}_{\mathrm{F}}, \boldsymbol{\beta}_{\mathrm{F}})
  &= \mathbf{W}_{\mathrm{F}}\, \mathrm{GAP}(\mathbf{H}_b) + \mathbf{b}_{\mathrm{F}},\\
\tilde{\mathbf{H}}_b &= (1 + \boldsymbol{\gamma}_{\mathrm{F}}) \odot \mathbf{H}_b + \boldsymbol{\beta}_{\mathrm{F}}
\end{aligned}
\label{eq:film}
\end{equation}
이다. 배율을 $1 + \boldsymbol{\gamma}_{\mathrm{F}}$ 로 두어 초기에는
항등 변조에서 출발한다. 국소 합성곱만으로는 ``지금 전황이 급박한가''와 같은 전역 정보가
모든 위치에 고르게 전달되지 않기 때문에 이 경로를 둔다.

각 방어정 $p$ 는 자신의 상태 벡터 $\mathbf{o}_p \in \mathbb{R}^8$(식~\eqref{eq:own})로부터
질의 벡터 $\mathbf{q}_p = \mathrm{MLP}_{8 \to 128 \to 32}(\mathbf{o}_p) \in \mathbb{R}^{32}$ 를
만들고, 픽셀 $i$ 의 특징 $\mathbf{f}_i \in \mathbb{R}^{32}$ 와의 내적으로 점수 $s_{p,i}$ 를
얻는다. 질의를 특징 공간에 맞추는 아핀 사영을 $\mathbf{W}_q \in \mathbb{R}^{d \times d}$,
$\mathbf{b}_q \in \mathbb{R}^{d}$ ($d = 32$)라 하면
\begin{equation}
s_{p,i} = \frac{\langle \mathbf{W}_q \mathbf{q}_p + \mathbf{b}_q,\, \mathbf{f}_i \rangle}{\sqrt{d}}, \qquad d = 32
\label{eq:score}
\end{equation}
이다. 이는 어텐션~\citep{vaswani2017attention}의 점수 계산과 같은 형식이며, 포인터
방식~\citep{vinyals2015pointer}으로 격자 위의 원소를 지목하는 구조에 해당한다. 특징 지도는
방어정마다 다르므로(방어정별 3채널), 관측 텐서는 방어정을 배치 축으로 쌓아 \textbf{결정당 한
번의 배치 호출}로 U-Net을 통과한다. 합성곱 연산량은 방어정 수에 비례하지만 \textbf{가중치는
모든 방어정이 공유하므로 모델 크기는 방어정 수에 불변}이다.

합성곱 블록은 $3\times3$ 합성곱, 그룹 정규화(8그룹)~\citep{wu2018groupnorm}, SiLU 의 순이다.
그룹 정규화는 표본 단위로 정규화하므로 배치 크기 1의 온보드 추론이 학습과 같은 연산이 된다.
배치 관점에서는 또한 정책이 조타 대신 경유점만 내어
선체 동역학~\citep{fossen2011marine}과 분리되고, 관측·행동이 격자 상대 좌표라 실해역
이식 시 축척만 재설정하면 된다. 다만 세 성질은 시뮬레이션 안에서만 확인되었다(\ref{sec:limitations}절).

\subsubsection{유효 마스크와 픽셀 지목}
\label{sec:validmask}

점수맵 전체가 후보가 되는 것은 아니다. 방어정 $p$ 의 유효 픽셀 집합 $\mathcal{V}_p$ 는 환상
요격 구역 안, 육지 밖, 배정 요격점의 반경 조건 안, 배정 무리의 방위 조건 안이라는 네
조건의 교집합이다. 픽셀 $i$ 의 중심을 $\mathbf{x}_i$, 육지 픽셀 집합을
$\mathcal{L}_{\mathrm{land}}$, 배정 요격점을 $\mathbf{I}_p$, 모선에서 본 요격점 방위와 픽셀
방위를 $\vartheta_p$, $\vartheta_i$, $[0^\circ, 180^\circ]$ 로 접은 각도 차이를 $|\cdot|_{\circ}$
라 하고, 요격 구역의 안·바깥 반경을 $r_{\min} = 400$, $r_{\max} = 4500\unit{m}$, 요격점 반경
조건을 $r_{\mathrm{g}} = 3000\unit{m}$, 방위 조건을 $\vartheta_{\mathrm{g}} = 60^\circ$ 로 두면
\begin{equation}
\begin{aligned}
\mathcal{V}_p = \big\{\, i :\; & r_{\min} \le \lVert \mathbf{x}_i - \mathbf{m} \rVert \le r_{\max},\;
  i \notin \mathcal{L}_{\mathrm{land}},\\
  & \lVert \mathbf{x}_i - \mathbf{I}_p \rVert \le r_{\mathrm{g}},\;
  |\vartheta_i - \vartheta_p|_{\circ} \le \vartheta_{\mathrm{g}} \,\big\}
\end{aligned}
\label{eq:validset}
\end{equation}
이다. 유효 픽셀이 하한 150개에 못 미치면 반경 조건부터 푼다. 방위 조건은 ``무엇을 막는
그물인가''를 정하므로 마지막까지 유지한다. 마스크는 로짓 단계에서 적용된다. 유효 집합
밖의 점수를 $-\infty$ 로 바꾼 마스크 점수 $\tilde{s}_{p,i}$ 를 소프트맥스에 넣어, 상태
$\mathbf{s}$ 에서 방어정 $p$ 가 픽셀 $i$ 를 고를 확률 $\pi_\theta(i \mid \mathbf{s}, p)$ 를
얻는다($\theta$: 정책 매개변수).
\begin{equationsmall}
\tilde{s}_{p,i} = \begin{cases} s_{p,i}, & i \in \mathcal{V}_p \\ -\infty, & i \notin \mathcal{V}_p, \end{cases}
\quad
\pi_\theta(i \mid \mathbf{s}, p) = \frac{e^{\tilde{s}_{p,i}}}{\sum_{j \in \mathcal{V}_p} e^{\tilde{s}_{p,j}}}
\label{eq:masked_softmax}
\end{equationsmall}
유효 집합 밖 픽셀의 확률은 정확히 0 이므로, 안전·임무 제약은 손실 항이 아니라 \emph{행동
분포에 직접} 부과되어 위반이 애초에 표현 불가능하다. 학습된 체크포인트로 잰 유효 비율은
격자의 평균 10.5\,\% 에 그쳐(Fig.~\ref{fig:validmask}(c)) 격자 대부분이 갈 수 없는 자리이며,
연속 좌표 회귀에는 이 구조를 표현할 수단이 없다.

\begin{figure*}[tbp]
\centering
\includegraphics[width=\textwidth]{fig9.png}
\caption{유효 마스크와 픽셀 지목(학습된 체크포인트의 실제 출력): (a) 유효 영역, (b) 점수맵과 선택된 픽셀 쌍,
         (c) 제약별 후보 수 감소(대표 사례에서 $2500 \to 1000 \to 979 \to 387 \to 299$ 칸,
         6{,}652회 결정 평균 유효 비율 10.5\,\%)}
\label{fig:validmask}
\end{figure*}

이산 격자의 또 다른 이득은 범주형 표본이 유효 영역을 벗어나지 않고 엔트로피가 유한하여
\ref{sec:training}절의 그룹 후보들이 실제로 서로 다른 칸을 밟는다는 점이다(가우시안 정책은
표준편차가 줄면 후보가 한 점으로 뭉친다). 한 칸 252\,m 의 해상도는 아래의 칸 안 연속
오프셋으로 보완한다.

각 방어정은 유효 영역 안에서 점수 상위 픽셀 $K=2$ 개를 \emph{순차적으로} 고른다. $k$ 번째
선택의 질의·후보 집합·선택 분포·선택된 픽셀을 $\mathbf{q}^{(k)}_p$, $\mathcal{V}^{(k)}_p$,
$\pi^{(k)}_p$, $i_k$ 라 하고($\mathbf{q}^{(0)}_p = \mathbf{q}_p$, $\mathcal{V}^{(0)}_p =
\mathcal{V}_p$), 두 점 사이 간격의 하한을 중복 배제 반경 $r_{\mathrm{dup}} = 200\unit{m}$, 상한을
그물벽 길이의 1.5배($1.5\,L_{\mathrm{net}} = 675\,$m)로 두면
\begin{equationfn}
\begin{aligned}
i_k &\sim \pi^{(k)}_p \equiv \pi_\theta(\cdot \mid \mathbf{q}^{(k)}_p, \mathcal{V}^{(k)}_p),\\
\mathbf{q}^{(k+1)}_p &= \mathbf{q}^{(k)}_p + \mathbf{f}_{i_k},\\
\mathcal{V}^{(k+1)}_p &= \big\{ i \in \mathcal{V}^{(k)}_p :
  r_{\mathrm{dup}} \le \lVert \mathbf{x}_i - \mathbf{x}_{i_k} \rVert \le 1.5\, L_{\mathrm{net}} \big\}
\end{aligned}
\label{eq:autoreg}
\end{equationfn}
이다. 먼저 고른 픽셀의 특징을 질의에 더해 두 번째를 고르므로 두 점이 서로를 알고 뽑히며,
셋째 줄의 하한은 벽이 서지 않을 만큼 가까운 두 점을, 상한은 그물 길이를 넘는 벽을 배제한다. 제약으로 후보가 모두 사라지면 제약을 완화하여 결정을 반환한다.

선택된 칸 안의 위치는 연속 오프셋으로 정한다. 첫 질의 $\mathbf{q}^{(0)}_p$ 와 선택 픽셀의
특징 $\mathbf{f}_{i_k}$ 로부터 MLP 로 평균 $\boldsymbol{\mu}_k$ 를 내고, 학습되는(로그 값으로
매개화한) 표준편차 $\boldsymbol{\sigma}_o$ 의 가우시안에서 정규화 오프셋 $\mathbf{u}_k$ 를
표집한 뒤, $[-1,1]$ 로 잘라 반픽셀 $\ell/2$ 를 곱해 픽셀 중심 $\mathbf{x}_{i_k}$ 에 더한 것이
$k$ 번째 목표점 $\mathbf{t}_k$ 이다.
\begin{equationsmall}
\begin{aligned}
\boldsymbol{\mu}_k &= \mathrm{MLP}_{2d \to d \to 2}\big([\mathbf{q}^{(0)}_p;\, \mathbf{f}_{i_k}]\big), \quad
\mathbf{u}_k \sim \mathcal{N}(\boldsymbol{\mu}_k, \boldsymbol{\sigma}_o^2),\\
\mathbf{t}_k &= \mathbf{x}_{i_k} + \tfrac{\ell}{2}\,\mathrm{clip}(\mathbf{u}_k, -1, 1)
\end{aligned}
\label{eq:offset}
\end{equationsmall}
즉 \emph{어느 칸인가}는 범주형으로, \emph{칸 안 어디인가}는 가우시안으로 결정된다. 한
방어정의 행동 $a_p = \{(i_k, \mathbf{u}_k)\}_{k=0}^{K-1}$ 의 로그확률은 두 분포의 합이며,
엔트로피 $\mathcal{H}(a_p)$ 도 같은 꼴로 정의된다.
\begin{equationsmall}
\log \pi_\theta(a_p) = \sum_{k=0}^{K-1} \Big[ \log \pi^{(k)}_p(i_k)
  + \log \mathcal{N}(\mathbf{u}_k; \boldsymbol{\mu}_k, \boldsymbol{\sigma}_o^2) \Big]
\label{eq:logp}
\end{equationsmall}
특징 조회에서 두 번째 점까지의 흐름을 Fig.~\ref{fig:picksteps}에 요약하였다.

\begin{figure*}[!tbp]
\centering
\includegraphics[width=0.9\textwidth]{fig10.png}
\caption{점수맵 위 픽셀 지목의 네 단계: 픽셀 특징, 유효 마스크, 첫 점 선택, 간격 제약 아래 둘째 점 선택}
\label{fig:picksteps}
\end{figure*}

선택된 두 좌표를 잇는 선분이 그물벽 경로가 된다. 방어정은 가까운 점까지 무전개로 이동한 뒤
다음 구간에서 그물을 전개한다. 운용 시에는 최대 점수를 취하고(greedy), 학습 시에는 확률
분포에서 표본을 뽑는다.

\textbf{정지 규칙.} 그물이 없거나 배정되지 않은 방어정은 정지시켜, 보상 신호 없이 평평해진
점수맵(최고 확률 0.02)이 항해 명령으로 새어 나가지 않게 한다.

\subsection{가치망 없는 다중 에이전트 그룹 상대 최적화}
\label{sec:training}

\subsubsection{그룹 상대 이득}

\begin{figure*}[tp]
\centering
\includegraphics[width=0.6\textwidth]{fig11.png}
\caption{가치망 없는 그룹 상대 최적화의 한 갱신: 후보 표집(위), 롤아웃 창 평가(가운데), 상대 비교와
         반사실적 크레딧(아래). 그림은 $G=4$, $H=4$ 로 줄여 그렸으며 실제 학습 설정은
         Table~\ref{tab:hparams}에 있다.}
\label{fig:grpocredit}
\end{figure*}

학습은 가치함수를 따로 두지 않는 그룹 상대 정책 최적화~\citep{shao2024deepseekmath}를 다중
방어정 상황에 맞게 변형하여 수행한다. 결정 시점의 상태를 $\mathbf{s}$, 방어정 $P$척의 결합
행동을 $\mathbf{a} = (a_1, \dots, a_P)$ 라 하자. 매 결정 시점에서 정책으로부터 $G$개의 결합
행동 후보를 표집하되, 후보 가운데 하나($k=0$)는 휴리스틱(\ref{sec:heuristic}절)이 내는
행동으로 고정하여 기준선으로 삼는다. $k$ 번째 후보를 $\mathbf{a}^{(k)}$, 휴리스틱이 상태
$\mathbf{s}$ 에서 내는 결합 행동을 $\mathbf{a}_{\mathrm{heur}}(\mathbf{s})$ 라 하면
\begin{equationfn}
\mathbf{a}^{(0)} = \mathbf{a}_{\mathrm{heur}}(\mathbf{s}), \qquad
\mathbf{a}^{(k)} \sim \pi_\theta(\cdot \mid \mathbf{s}), \quad k = 1, \dots, G-1
\label{eq:candidates}
\end{equationfn}
각 후보를 시뮬레이터에서 공통 구간 $H = 180\,\mathrm{step}$ 만큼 실제로 진행시켜 창 보상
$R_k = R(\mathbf{s}, \mathbf{a}^{(k)}; H)$ 를 얻은 뒤(\ref{sec:reward}절), 후보들 사이의
상대 비교로 이득 $A_k$ 를 산정한다. 기준은 휴리스틱 후보의 보상 $R_0$ 이고, 그룹 표준편차로
나누되 0 나눗셈을 막는 $\epsilon$ 을 더한다.
\begin{equation}
A_k = \frac{R_k - R_0}{\mathrm{std}(\{R_j\}_{j=0}^{G-1}) + \epsilon}, \qquad \epsilon = 10^{-6}
\label{eq:adv}
\end{equation}
$R_0 = R_{\mathrm{heur}}$ 이므로 휴리스틱과 같은 성과를 내는 후보의 이득은 0 이고, 휴리스틱을
넘어서는 후보만 양의 이득을 받는다. 그룹 표준편차로 나누므로 이득은 크기가 없는 양이며,
그룹 안의 순위 정보만 남는다. 어떤 선택이 더 나았는지를 실제 결과로 직접 비교하므로
가치망과 그 추정 편향·학습 불안정이 없다.

한 갱신은 $N$개의 병렬 전장(월드)에서 동시에 수행된다. 월드 $n$ 의 후보 보상이 모두 같으면
이득이 정의되지 않으므로, 유효 마스크 $w_n = \mathbb{1}[\mathrm{std}(\{R_{n,k}\}_k) > \epsilon_{\mathrm{var}}]$
($\epsilon_{\mathrm{var}} = 10^{-6}$)로 그런 월드를 갱신에서 제외한다. 월드 $n$ 의 상태를
$\mathbf{s}_n$, 후보 $k$ 의 이득을 $A_{n,k}$, 그 후보에서 방어정 $p$ 의 행동을 $a^{(k)}_{n,p}$
라 하면, 정책경사 손실 $\mathcal{L}_{\mathrm{PG}}$ 는 유효 월드의 모든 후보·모든 방어정에
대해 이득으로 가중한 로그확률(식~\eqref{eq:logp})의 음수이다.
\begin{equationfn}
\mathcal{L}_{\mathrm{PG}}(\theta) = -\frac{\sum_{n} w_n \sum_{k=0}^{G-1} A_{n,k} \sum_{p=1}^{P} \log \pi_\theta\big(a^{(k)}_{n,p} \mid \mathbf{s}_n\big)}{G \sum_n w_n}
\label{eq:lpg}
\end{equationfn}
여기에 계수 $\lambda_{\mathrm{H}}(t)$(\ref{sec:schedule}절의 감쇠 스케줄)의 엔트로피
정칙화를 더한 총손실을 최소화한다.
\begin{equationsmall}
\mathcal{L}(\theta) = \mathcal{L}_{\mathrm{PG}}(\theta)
  - \lambda_{\mathrm{H}}(t)\, \frac{1}{NP} \sum_{k=0}^{G-1} \sum_{n=1}^{N} \sum_{p=1}^{P} \mathcal{H}\big(a^{(k)}_{n,p}\big)
\label{eq:loss}
\end{equationsmall}
엔트로피는 $G$개 후보 모두에 대해 합하므로 유효 계수는 $G \lambda_{\mathrm{H}}(t)$ 이다.
기울기 노름을 클립한 뒤
갱신하며, 관련 상수는 Table~\ref{tab:hparams}에 모았다. 이 절차와 다음 소절의 크레딧
배분을 Fig.~\ref{fig:grpocredit}에 정리하였다.

\subsubsection{반사실적 크레딧 배분: 비교 변형}

식~\eqref{eq:lpg}는 팀 보상의 이득을 모든 방어정에 같은 값으로 곱하므로 개별 기여가 가려질
수 있다. 이를 분리하는 변형으로, 기준 표본 $\mathbf{a}^{\mathrm{base}} \sim
\pi_\theta(\cdot \mid \mathbf{s})$ 에서 방어정 $p$ 의 행동만 바꾼 혼합 후보 $G$개(후보 0 은
휴리스틱)를 같은 구간 $H$ 만큼 재진행하여 창 보상 $R^{(p)}_k$ 를 얻고, 방어정 $p$ 만의 이득
\begin{equation}
A^{(p)}_k = \frac{R^{(p)}_k - R^{(p)}_0}{\mathrm{std}(\{R^{(p)}_j\}_{j=0}^{G-1}) + \epsilon}
\label{eq:cfadv}
\end{equation}
을 방어정 $p$ 의 로그확률에만 곱한다. 나머지 방어정이 고정되어 있으므로 이 차이는 $p$ 의
선택에서만 온다. COMA~\citep{foerster2018coma}의 반사실적 기준선과 같은 착상이되 학습된
비평가 대신 시뮬레이터를 실제로 재진행한다. 본 논문이 평가하는 정책은 \emph{팀 이득
손실}(식~\eqref{eq:lpg}--\eqref{eq:loss})로 학습한 것이며, 같은 액터 위에서 갱신 규칙만
MAPPO 로 바꾼 기준선과의 비교는 \ref{sec:res_algo}절에 둔다.

\subsubsection{보상 설계}
\label{sec:reward}

한 후보의 창 보상 $R(\mathbf{s}, \mathbf{a}; H)$ 는 결정 시점부터 $H$\,step 동안 일어난
일을 하나의 스칼라로 합산한 팀 공유 보상이며, 이벤트 항 $R_\mathrm{event}$, 창의 시작·끝
상태 $\mathbf{s}_t$, $\mathbf{s}_{t+H}$ 에서 잰 퍼텐셜 $\Phi$ 의 할인($\gamma$) 차분, 조밀 항
$R_\mathrm{dense}$ 로 구성된다. 세 항의 계수는 Table~\ref{tab:reward_coef}에 모았다.
\begin{equation}
R = R_\mathrm{event} + \big[\gamma\,\Phi(\mathbf{s}_{t+H}) - \Phi(\mathbf{s}_t)\big] + R_\mathrm{dense}
\label{eq:reward}
\end{equation}

\begin{table}[tp]
\centering
\caption{보상 계수. 이벤트 항은 사건 1건당 가중, 퍼텐셜·조밀 항은 정규화된 수량에 곱하는
         가중과 정규화 상수다.}
\label{tab:reward_coef}
\footnotesize
\renewcommand{\arraystretch}{1.15}
\begin{tabular}{@{}llr@{\hspace{2mm}}l@{}}
\toprule
항 & 기호 & 값 & 의미 \\
\midrule
\multirow{6}{*}{이벤트}
 & $r_{\mathrm{cap}}$  & $+13.5$ & 포획 1건 \\
 & $r_{\mathrm{brc}}$  & $-8.0$  & 돌파 1건 \\
 & $r_{\mathrm{col}}$  & $-30$   & 충돌·육지 충돌 1건 \\
 & $r_{\mathrm{tng}}$  & $-13$   & 그물 접촉 1건 \\
 & $r_{\mathrm{wipe}}$ & $+8$    & 창 안 종료 시 전멸 \\
 & $r_{\mathrm{surv}}$ & $-5$    & 창 안 종료 시 잔존 적 1척당 \\
\midrule
\multirow{3}{*}{퍼텐셜}
 & $w_{\mathrm{threat}}$ & $3.0$  & 모선 근접도 합의 가중 \\
 & $R_w$                 & $6300\unit{m}$ & 근접도 정규화(교전 해역 반폭) \\
 & $\gamma$              & $0.99$ & 할인 \\
\midrule
\multirow{10}{*}{조밀}
 & $w_{\mathrm{path}}$ & $0.5$  & 이동거리 비용 \\
 & $w_{\mathrm{net}}$  & $0.3$  & 그물 소모 비용 \\
 & $w_{\mathrm{turn}}$ & $0.3$  & 선회량 비용 \\
 & $w_{\mathrm{cons}}$ & $0.3$  & 목표점 일관성 비용 \\
 & $w_{\mathrm{prox}}$ & $0.75$ & 모선·타 방어정 근접 비용 \\
 & $w_{\mathrm{cov}}$  & $0.1$  & 커버리지 보너스 \\
 & $w_{\mathrm{ray}}$  & $1.6$  & 배율판 보너스($\mathrm{Ray} \in [0, 2.6]$) \\
 & $R_{\mathrm{m}}$      & $800\unit{m}$ & 모선 안전 반경 \\
 & $\varrho_{\mathrm{m}}$ & $500\unit{m}$ & 모선 근접 영향 반경 \\
 & $\varrho_{\mathrm{a}}$ & $600\unit{m}$ & 타 방어정 근접 영향 반경 \\
\bottomrule
\end{tabular}
\end{table}

이벤트 항은 창 안의 사건 수에 고정 가중을 곱한 합이다. 사건 수는 \ref{sec:metrics}절의
$N_{\mathrm{cap}}$, $N_{\mathrm{brc}}$, $N_{\mathrm{col}}$, $N_{\mathrm{tng}}$ 에 육지 충돌
$N_{\mathrm{land}}$ 를 더한 것이고, 창 안에서 에피소드가 끝나면($\mathbb{1}_{\mathrm{end}}$)
전멸 $r_{\mathrm{wipe}}$ 와 잔존 적 1척당 $r_{\mathrm{surv}}$ 를 더한다.
\begin{equationsmall}
\begin{aligned}
R_\mathrm{event} = \;& r_{\mathrm{cap}} N_{\mathrm{cap}} + r_{\mathrm{brc}} N_{\mathrm{brc}}
  + r_{\mathrm{col}} \big(N_{\mathrm{col}} + N_{\mathrm{land}}\big)\\
  & + r_{\mathrm{tng}} N_{\mathrm{tng}}
  + \mathbb{1}_{\mathrm{end}}\big(r_{\mathrm{wipe}}\,\mathbb{1}[\mathcal{A} = \varnothing] + r_{\mathrm{surv}}\,|\mathcal{A}|\big)
\end{aligned}
\label{eq:revent}
\end{equationsmall}
퍼텐셜 $\Phi$ 는 생존 적선 집합 $\mathcal{A}$ 의 모선 근접도(적선까지 거리를 교전 해역
반폭 $R_w$ 로 나눈 값을 1에서 뺀 것)의 합에 가중 $w_\mathrm{threat}$ 을 곱해 부호를 바꾼
것으로, 할인 $\gamma$ 의 차분 형태로 주어지므로 최적 정책을 보존한다~\citep{ng1999shaping}.
\begin{equation}
\Phi(\mathbf{s}) = -w_\mathrm{threat} \sum_{m \in \mathcal{A}}
          \mathrm{clip}\!\left(1 - \frac{\lVert \mathbf{e}_m - \mathbf{m} \rVert}{R_w},\, 0,\, 1\right)
\label{eq:potential}
\end{equation}
적이 포획되어 $\mathcal{A}$ 에서 빠지거나
멀어지면 $\Phi$ 가 오른다.

조밀 항 $R_\mathrm{dense}$ 는 창 안의 기동 품질에 대한 비용·보너스의 합이다. 창 길이 $H$,
방어정 속력 $v_a$, 창 안 총 이동거리 $D$, 소모 그물 수 $N_{\mathrm{net}}$ 에 더해, 방어정 $p$
의 창 안 누적 선회량을 최대 가능 선회량 $H\dot\psi_{\max}$ 로 나눈 값을 $S_p \in [0,1]$,
직전 결정 대비 목표점 이동을 $L_{\mathrm{net}}$ 으로 나눈 값을 $\Xi_p$, 모선·타 방어정에 대한
근접도를 $\rho^{\mathrm{m}}_p$, $\rho^{\mathrm{a}}_p$(식~\eqref{eq:prox}), 설치된 그물이 적의
접근로를 막는 비율을 $\mathrm{Cov}$(식~\eqref{eq:coverage}), 계획된 벽까지 포함한 그
배율판을 $\mathrm{Ray}$ 라 하면
\begin{equationfn}
\begin{aligned}
R_\mathrm{dense} = \;& -\,w_{\mathrm{path}} \frac{D}{H v_a P} - w_{\mathrm{net}} N_{\mathrm{net}}\\
  & - \frac{w_{\mathrm{turn}}}{P}\sum_{p=1}^{P} S_p
    - \frac{w_{\mathrm{cons}}}{P}\sum_{p=1}^{P} \Xi_p\\
  & - w_{\mathrm{prox}} \sum_{p=1}^{P}\big(\rho_{p}^{\mathrm{m}\,2} + \rho_{p}^{\mathrm{a}\,2}\big)
    + w_{\mathrm{cov}}\, \mathrm{Cov} + w_{\mathrm{ray}}\, \mathrm{Ray}
\end{aligned}
\label{eq:rdense}
\end{equationfn}
이다. 근접도는 창 안에서 방어정 $p$ 가 모선·타 방어정에 가장 가까웠던 거리 $d^{\mathrm{m}}_p$,
$d^{\mathrm{a}}_p$ 를 모선 안전 반경 $R_{\mathrm{m}}$ 와 영향 반경 $\varrho_{\mathrm{m}}$,
$\varrho_{\mathrm{a}}$ 로 정규화한 것이며, 조밀 항에서는 제곱하여 붙을수록 페널티가 급격히
커지게 한다.
\begin{equationsmall}
\rho^{\mathrm{m}}_p = \mathrm{clip}\!\Big(\frac{R_{\mathrm{m}} - d^{\mathrm{m}}_p}{\varrho_{\mathrm{m}}}, 0, 1\Big), \quad
\rho^{\mathrm{a}}_p = \mathrm{clip}\!\Big(\frac{\varrho_{\mathrm{a}} - d^{\mathrm{a}}_p}{\varrho_{\mathrm{a}}}, 0, 1\Big)
\label{eq:prox}
\end{equationsmall}
커버리지 $\mathrm{Cov}$ 는 적에서 모선으로 향하는 직선(ray)이 설치된 그물 영역 $\mathcal{W}$
에 막힌 생존 적의 비율로, 선분 위 표본점 $\lambda \in \Lambda$($[0.05, 0.95]$ 를 16등분)에서
검사한다. 창 안에 포획이 없을 때도 후보를 변별하는 선행 지표다.
\begin{equationsmall}
\mathrm{Cov} = \frac{1}{|\mathcal{A}|} \sum_{m \in \mathcal{A}}
  \mathbb{1}\big[\exists\, \lambda \in \Lambda :\; \mathbf{e}_m + \lambda(\mathbf{m} - \mathbf{e}_m) \in \mathcal{W}\big]
\label{eq:coverage}
\end{equationsmall}
$\mathrm{Ray}$ 는 같은 검사를 계획된 그물벽까지 포함하되 적에 가까운 차단을 우대한
배율판이다. 퍼텐셜 항은 이벤트 항의 희소성을 메운다. 216개 결정의 실측에서 포획은
평가 창의 46.8\,\% 에서만 발생했고, 나머지 창에서는 퍼텐셜 항이 후보 순위를 가르는 주된
신호였다.

\subsubsection{초기화, 스케줄, 수렴 진단}
\label{sec:trainproc}

\paragraph{행동 복제 초기화}
\label{sec:bc}
그룹 상대 이득은 후보들이 갈릴 때만 신호를 주므로, 휴리스틱 행동(\ref{sec:heuristic}절)을
목표로 300 step 의 행동 복제로 정책을 부팅한다. 목표는 원-핫이 아니라 휴리스틱 픽셀
$i^{*}_k$ 를 중심으로 폭 $\sigma_{\mathrm{BC}} = 0.8$\,px($\approx 202\unit{m}$)의 가우시안을
유효 픽셀 위에서 정규화한 가중 $\omega_k(i)$ 이고(인접 픽셀은 사실상 같은 행동이라 원-핫은
기울기가 지나치게 희소하다), 손실 $\mathcal{L}_{\mathrm{BC}}$ 는 이 가중을 건 교차 엔트로피다.
\begin{equation}
\omega_k(i) = \frac{\mathbb{1}[i \in \mathcal{V}^{(k)}_p]\,
  \exp\!\big(-\lVert \mathbf{x}_i - \mathbf{x}_{i^{*}_k} \rVert^2 / 2\sigma_{\mathrm{BC}}^2\big)}
  {\sum_{j \in \mathcal{V}^{(k)}_p} \exp\!\big(-\lVert \mathbf{x}_j - \mathbf{x}_{i^{*}_k} \rVert^2 / 2\sigma_{\mathrm{BC}}^2\big)}
\label{eq:bcweight}
\end{equation}
\begin{equation}
\mathcal{L}_{\mathrm{BC}}(\theta) = -\frac{1}{P}\sum_{p=1}^{P} \sum_{k=0}^{K-1} \sum_{i \in \mathcal{V}^{(k)}_p}
  \omega_k(i)\, \log \pi^{(k)}_p(i)
\label{eq:lbc}
\end{equation}
두 번째 픽셀의 질의는 휴리스틱의 첫 픽셀로 갱신한다(teacher forcing).

\paragraph{학습률·엔트로피 계수 스케줄}
\label{sec:schedule}
옵티마이저는 Adam 이다. 학습률은 $\eta_0 = 10^{-4}$ 에서 $0.05\,\eta_0$ 까지 단일 주기
코사인 어닐링~\citep{loshchilov2017sgdr}($T = 800$ 업데이트)으로, 식~\eqref{eq:loss}의
엔트로피 계수는 $\lambda_0 = 0.01$ 에서 $0.1\,\lambda_0$ 까지 같은 예산 축 위에서 선형으로
줄인다. 두 계수를 같은 축에서 줄이면 분포가 뾰족해지는 시점과 보폭이 작아지는 시점이
맞물린다. 채택 가중치(600 업데이트, \ref{sec:protocol}절)는 $\eta \approx 0.19\,\eta_0$,
$\lambda \approx 0.33\,\lambda_0$ 인 감쇠 후반에 놓인다.

\begin{figure*}[tp]
\centering
\includegraphics[width=\textwidth]{fig12.png}
\caption{수렴 진단(회색 시드별, 파란 실선 시드 평균): (a) greedy 포획률, (b) 정책 엔트로피, (c) 근사 KL 발산,
         (d) 파라미터 이동}
\label{fig:convergence}
\end{figure*}

\paragraph{수렴 진단}
\label{sec:convergence}
채택 런은 20 업데이트마다의 요약만 남겼으므로, 같은 설정을 시드 3개로 다시 돌려 매
업데이트마다 수렴 진단을 기록하였다(Fig.~\ref{fig:convergence}; (a)의 빨간 파선은 채택 런,
초록 파선은 휴리스틱 기준선). 진단량은 갱신 전후 정책 사이의 근사 KL(같은 관측·같은 표본
후보에 대한 확률비의 $k_3$ 추정량)과 행동복제 초기점으로부터의 상대 파라미터 이동이다.
세 시드에서 근사 KL은 처음 100 업데이트 평균 $1.4\times10^{-2}$ 에서 마지막 100 업데이트
$1.5\times10^{-4}$ 로 두 자릿수 줄고, 파라미터 이동은 뒤 절반에서 포화하며, 정책 엔트로피는
$2.42$ 에서 $2.06$\,nat 로 내려가 400 업데이트 이후 평탄하다. greedy 포획률은 세 시드 모두
800 업데이트에서 기준선 $0.749$ 위에 있고 채택 런의 평가점도 그 범위 안에 든다.

\subsection{언어모델 기반 전략 계층}
\label{sec:commander}
\label{sec:schema}

전략 계층은 Fig.~\ref{fig:architecture}의 상단에 해당한다. 입력은 전장 상태·임무 목표·교전
규칙·자산 상태이고, 출력은 무리별 배정과 대기(HOLD)이며, 그 배정이 기동 계층으로 전달되어
각 방어정의 그물벽 위치 결정을 조건짓는다(실행 화면은 Fig.~\ref{fig:runshot}).

\begin{figure*}[tp]
\centering
\includegraphics[width=0.9\textwidth]{fig13.png}
\caption{실행 화면(양동 포메이션, Gemini 지휘관): 전장(왼쪽)과 지휘관의 배정·근거(오른쪽)}
\label{fig:runshot}
\end{figure*}

언어모델에는 좌표 원시 데이터 대신 \emph{기하량이 미리 계산된} 전장 상태를 텍스트로
전달한다. 각 적 무리의 요격점까지의 이동거리와 선회량, 모선 기준 방위, 설치된 그물에
의한 접근로 차단 여부가 포함된다.

출력은 자유 문장이 아니라 다음 자료형에서 생성한 JSON 스키마로
강제한다~\citep{openai2026function}. 자유 문장은 하위 계층의 파싱 실패로 그 주기의 배정을
통째로 잃을 수 있는 반면, 스키마는 디코딩 단계에서 문법으로 작용하여 형식에 맞지 않는
토큰 자체가 생성되지 않게 한다.

% 좁은 단에 맞춰 줄을 미리 끊는다. 자동 줄바꿈에 맡기면 들여쓰기가 무너져
% 자료형 구조가 안 보인다.
\begin{quote}\ttfamily\scriptsize\setlength{\parindent}{0pt}
CommanderPlan \{\\
\hspace*{1.2em} rationale: str\\
\hspace*{1.2em} deployments: List[\{\\
\hspace*{3.0em} cluster\_id: int,\\
\hspace*{3.0em} ally\_ids: List[int] \}]\\
\hspace*{1.2em} hold\_ships: List[int]\\
\}
\end{quote}

언어모델이 정하는 것은 \emph{어느 배가 어느 무리를 맡고 누가 대기하는가}뿐이며, 그물
전개 여부·벽의 방위·교전 반경과 요격점 기하는 기동 계층이 정한다. 배정은 코드가 보정하지
않고 그대로 쓰인다. 비워 둔 \texttt{ally\_ids}는 그 무리를 무방비로 두고, 지목되지 않은
방어정은 대기한다. 따라서 \ref{sec:results}절의 조건 간 차이는 배정 판단 자체에 귀속된다.
온도는 0으로 두어 같은 전장에서 같은 배정이 나오게 하며, 이는 \ref{sec:paired}절의 시드
대응 설계의 전제이기도 하다.

한 무리에 두 척 이상이 배정되면 \ref{sec:heuristic}절의 잉여 배정과 같은 반경 층
분리($\lambda_{\mathrm{L}} = 0.18$)와 방위 분리($\pm 40^\circ$)가 작동하되, 요격점이 요격
구역을 벗어나 유효 마스크가 붕괴하지 않도록 반경을 구역 안으로 자른다. 배정이 바뀌면
요격점과 함께 기동 계층의 유효 마스크(\ref{sec:validmask}절)도 바뀌므로, 배정은 기동의
선행 조건이다.
